\subsection{ROOT SYSTEMS AND GROUPS GENERATED BY REFLECTIONS}

PROPOSITION 8. Let F be a real Hilbert space of finite dimension l. Let H be the set of affine hyperplanes of F and G the group generated by the orthogonal reflections \(s_{H}\) with respect to the hyperplanes \(H \in H\) . Assume that the conditions of Chap. V, § 3 are satisfied (i.e.~that \(g(H) \in H\) for all \(H \in H\) and \(g \in G\) , and that G acts properly on F). Assume also that 0 is a special point for G and that the group T of translations belonging to G is of rank l. Then there exists a unique reduced root system R in \(V = F^{*}\) such that the canonical isomorphism from F to \(V^{*}\) transforms G to the affine Weyl group \(W_{a}\) of R.

We remark first of all that the assumption on T implies that G is essential: otherwise, the affine space F would decompose into a product \(F_{0} \times F_{1}\) , with \(\dim F_{1} < l\) , the group G being identified with a group of displacements acting properly on \(F_{1}\) (Chap. V, § 3, no. 8, Prop. 6), and T would not be of rank l.

Let \(\mathfrak{H}_0\) be the set of \(\mathbf{H} \in \mathfrak{H}\) such that \(0 \in \mathbf{H}\) . For \(\mathbf{H} \in \mathfrak{H}_0\) , let \(\mathfrak{H}_{\mathbf{H}}\) be the set of elements of \(\mathfrak{H}\) parallel to \(\mathbf{H}\) . Since 0 is a special point, \(\mathfrak{H}\) is the union of the \(\mathfrak{H}_{\mathbf{H}}\) for \(\mathbf{H} \in \mathfrak{H}_0\) . Since \(\mathbf{T}\) is of rank \(l\) , there exists a \(v \in \mathbf{F}\) such that the translation by the vector \(v\) belongs to \(\mathbf{T}\) and \(v \notin \mathbf{H}\) . The hyperplanes \(\mathbf{H} + kv\) for \(k \in \mathbf{Z}\) are pairwise distinct and belong to \(\mathfrak{H}_{\mathbf{H}}\) . Now let \(a\) be a unit vector of \(F\) orthogonal to \(\mathbf{H}\) : then \(\mathbf{H} + (v|a)a \in \mathfrak{H}_{\mathbf{H}}\) and since \(\mathfrak{H}\) is locally finite (Chap. V, § 1, no. 1, Lemma 1), there is a smallest real number \(\lambda > 0\) such that \(\mathbf{H} + \lambda a \in \mathfrak{H}\) . We are going to show that \(\mathfrak{H}_{\mathbf{H}}\) is the set of hyperplanes \(\mathbf{H} + k\lambda a\) for \(k \in \mathbf{Z}\) . Indeed,

\[
\mathrm{H} ^ {\prime} = \mathrm{H} + \lambda a \in \mathfrak {H} _ {\mathrm{H}}
\]

and the element \(s_{\mathrm{H}'} \circ s_{\mathrm{H}}\) of \(\mathbf{G}\) is the translation by the vector \(2\lambda a\) (Chap. V, § 2, no. 4, Prop. 5). Consequently, \(\mathrm{H} + 2n\lambda a = (s_{\mathrm{H}'}s_{\mathrm{H}})^n(\mathrm{H})\) and \(\mathrm{H} + (2n + 1)\lambda a = (s_{\mathrm{H}'}s_{\mathrm{H}})^n(\mathrm{H}')\) belongs to \(\mathfrak{H}_{\mathrm{H}}\) . On the other hand, if \(\mathrm{L} \in \mathfrak{H}_{\mathrm{H}}\) , there exists \(\xi \in \mathbf{R}\) such that \(\mathrm{L} = \mathrm{H} + \xi \lambda a\) and there exists an integer \(n\) such that

\[
\text { either } 2 n <   \xi \leqslant 2 n + 1, \text { or } 2 n - 1 <   \xi \leqslant 2 n.
\]

In the first case, \((s_{\mathrm{H}}s_{\mathrm{H}^{\prime}})^{n}(\mathrm{L}) = \mathrm{H} + (\xi -2n)\lambda a\) with

\[
0 <   (\xi - 2 n) \lambda \leqslant \lambda
\]

and the definition of \(\lambda\) implies that \(\xi = 2n + 1\) ; in the second case,

\[
s _ {\mathrm{H}} (s _ {\mathrm{H}} s _ {\mathrm {H^ {\prime}}}) ^ {n} (\mathrm{L}) = \mathrm{H} + (2 n - \xi) \lambda a \text {with} 0 \leqslant (2 n - \xi) \lambda <   \lambda
\]

and the definition of \(\lambda\) implies that \(\xi = 2n\) .

It follows that if \(\alpha_{\mathrm{H}}\) is the linear form on \(F\) such that

\[
\mathrm{H} ^ {\prime} = \left\{x \in \mathrm{F} \mid \langle \alpha_ {\mathrm{H}}, x \rangle = 1 \right\},
\]

the set \(\mathfrak{H}_{\mathrm{H}}\) is the set of hyperplanes \(\mathrm{L}_{\alpha_{\mathrm{H}},k} = \{x\in \mathrm{F}\mid \langle \alpha_{\mathrm{H}},x\rangle = k\}\) for \(k\in \mathbf{Z}\) , and \(\alpha_{\mathrm{H}}\) and \(-\alpha_{\mathrm{H}}\) are the only linear forms with this property.

Consequently, the proposition will be proved if we show that the set R of elements of V of the form \(\pm\alpha_{H}\) is a reduced root system in V.

\begin{enumerate}
\def\labelenumi{\alph{enumi})}
\item
  We prove condition \((\mathrm{RS}_{\mathrm{I}})\) : it is clear that R is finite (since \(H_{0}\) is finite) and does not contain 0. Moreover, R generates V. Indeed, if \(x \in F\) is orthogonal to R, then \(x \in H\) for all \(H \in H_{0}\) and the translation by the vector x commutes with every element of G. Since G is essential, this implies that x = 0.
\item
  We prove \((\mathrm{RS}_{\mathrm{II}})\) . For \(v \in V\) and \(r \in R\) , put \(L_{v,r} = \{x \in F | \langle v, x \rangle = r\}\) as above; if \(\alpha \in R\) , put \(H_{\alpha} = L_{\alpha,0}\) , and let \(s_{\alpha}\) be the transpose of \(s_{H_{\alpha}}\) . There exists a unique element \(\alpha' \in F\) orthogonal to \(H_{\alpha}\) and such that \(\langle \alpha', \alpha \rangle = 2\) . Then \(s_{H_{\alpha}} = s_{\alpha', \alpha}\) and \(s_{\alpha} = s_{\alpha, \alpha'}\) . For \(\beta \in R\) ,
\end{enumerate}

\[
\mathrm{L} _ {s _ {\alpha} (\beta), 1} = s _ {\mathrm{H} _ {\alpha}} \left(\mathrm{L} _ {\beta , 1}\right) \in \mathfrak {H}
\]

and there exist \(\gamma \in \mathbb{R}\) and \(n\in \mathbf{N}^*\) such that \(\mathrm{L}_{s_{\alpha}(\beta),1} = \mathrm{L}_{\gamma ,n}\) . Then

\[
s _ {\mathrm{H} _ {\alpha}} (\mathrm{L} _ {\gamma , 1}) = \mathrm{L} _ {\beta , 1 / n}
\]

and so \(1 / n\in \mathbf{Z}\) . Thus \(n = 1\) and \(s_{\alpha}(\beta) = \gamma \in \mathbb{R}\) . This proves \((\mathrm{RS}_{\mathrm{II}})\) .

\begin{enumerate}
\def\labelenumi{\alph{enumi})}
\setcounter{enumi}{2}
\tightlist
\item
  We prove \((\mathrm{RS}_{\mathrm{III}})\) . Let \(\alpha \in R\) and set \(H'_{\alpha} = L_{\alpha,1}\) . Then
\end{enumerate}

\[
\mathrm{H} _ {\alpha} ^ {\prime} = \mathrm{H} _ {\alpha} + (1 / 2) \alpha^ {\prime};
\]

since the translation \(t(\alpha^{\prime})\) by the vector \(\alpha^{\prime}\) is the product \(s_{\mathrm{H}_{\alpha}^{\prime}}s_{\mathrm{H}_{\alpha}}\) (Chap. V, § 2, no. 4, Prop. 5), it belongs to T and \(\alpha^{\prime} = t(\alpha^{\prime})(0)\) is a special point for G. Consequently, for all \(\beta \in \mathbb{R}\) , there exists a hyperplane \(\mathrm{L}_{\beta,k}\) passing through \(\alpha^{\prime}\) , with \(k\) an integer, which shows that \(\langle \beta ,\alpha^{\prime}\rangle \in \mathbf{Z}\) , and proves (RSIII).

\begin{enumerate}
\def\labelenumi{\alph{enumi})}
\setcounter{enumi}{3}
\tightlist
\item
  Finally, it is clear that R is reduced, for if \(H, H' \in H_{0}, H \neq H'\) , the linear forms \(\alpha_{H}\) and \(\alpha_{H'}\) are not proportional.
\end{enumerate}

Remark. 1) The assumption that T is of rank l is satisfied in particular when G is irreducible and infinite. Indeed, the vector space generated by vectors corresponding to the translations in T is invariant under the canonical image of G in the linear group of F. It is different from \(\{0\}\) if G is infinite and is thus equal to the whole of F if G is infinite and irreducible.

A finite group generated by reflections is not always the Weyl group of a root system. More precisely:

PROPOSITION 9. Let V be a real vector space of finite dimension l, and let G be a finite subgroup of \(\mathbf{GL}(V)\) , generated by reflections and essential. Give V a scalar product invariant under G. The following conditions are equivalent:

\begin{enumerate}
\def\labelenumi{(\roman{enumi})}
\item
  There exists a discrete subgroup of \(\mathrm{V}\) of rank \(l\) that is stable under \(\mathrm{G}\) .
\item
  There exists a \(\mathbf{Q}\) -structure on V (Algebra, Chap. II, § 8, no. 1, Def. 1) invariant under \(\mathbf{G}\) .
\item
  There exists a root system in V whose Weyl group is G.
\item
  There exists a discrete group \(G'\) of displacements of V, acting properly on V, and generated by reflections, such that \(G'\) is the semi-direct product of G and a group of translations of rank l.
\item
  \(\Longrightarrow\) (i): let \(V' \subseteq V\) be a Q-structure on V invariant under G. Let A be a finite subset of \(V'\) generating the Q-vector space \(V'\) . Replacing A by \(\bigcup_{s \in G} s(A)\) , we can assume that A is stable under G. Let B be the subgroup of V generated by A. Then B is stable under G, of finite type and torsion-free, so has a basis over Z which is both a basis of \(V'\) over Q and a basis of V over R.
\item
  \(\Longrightarrow\) (ii): this follows from Prop. 1 of § 1, no. 1, for example.
\item
  \(\Longrightarrow\) (iii): let \(G'\) be a group satisfying condition (iv). The group of translations of \(G'\) is of rank l, and 0 is a special point for \(G'\) by Prop. 9 of Chap. V, § 3, no. 10. Prop. 6 shows that there exists a reduced root system \(R_{0}\) in \(V^{*}\) such that \(G'\) is identified with \(W_{a}(R_{0})\) ; the group G is then the Weyl group of the inverse root system of \(R_{0}\) .
\item
  \(\Longrightarrow\) (iv): assume that G leaves stable a discrete subgroup M of V, of rank l. For any reflection \(s \in G\) , \(s(x) - x \in M\) for all \(x \in M\) , so the line \(D_{s}\) orthogonal to \(H_{s}\) meets M; let \(\alpha_{s}, -\alpha_{s}\) be the generators of the cyclic group \(D_{s} \cap M\) ; the set A of the \(\alpha_{s}\) and \(-\alpha_{s}\) is stable under G, hence generates a subgroup \(M'\) of M stable under G; the discrete group \(M'\) is of rank l because G is essential. Let \(G'\) be the group of affine transformations of V that is the semi-direct product of G and the group of translations whose vectors belong to \(M'\) . Let \(G_{1}'\) be the subgroup of \(G'\) generated by the reflections of \(G'\) . We shall show that \(G_{1}' = G'\) , which will complete the proof. First, \(G_{1}' \supseteq G\) since G is generated by reflections. On the other hand, for any reflection s of G, let \(t_{s}\) be the translation with vector \(\alpha_{s}\) . The transformation \(s \circ t_{s}\) is a reflection, and \(s \circ t_{s} \in G'\) ; thus \(t_{s}\) is a product of two reflections of \(G'\) ; this being true for every reflection s of G, the translations whose vector belongs to \(M'\) are all in \(G_{1}'\) .
\end{enumerate}

DEFINITION 2. A group G satisfying the equivalent conditions of Prop. 9 is called a crystallographic group.

Remark. 2) Let G be a finite group generated by reflections and essential. Then G is crystallographic if and only if every element of its Coxeter matrix is one of the integers 1, 2, 3, 4, 6. Indeed, this condition is necessary by Remark 3) of § 1, no. 5. The fact that it is sufficient will follow from the classification of finite Coxeter groups given in § 4 (for a direct proof, see Chap. V, § 4, Exerc. 6).

